\subsection{泛映射的存在性}

一个泛映射问题不一定有解（习题1）。然而，我们将证明以下条件蕴含解的存在性：

\((\mathbf{CU}_{\mathbf{I}})\) 在每一族 \(\Sigma\) -集合的乘积上，存在一个物种 \(\Sigma\) 的乘积结构（§2，第4节）。

\((\mathrm{CU}_{\mathrm{II}})\) 让 \((\mathbf{F}_i)_{i\in \mathbf{I}}\) 是 \(\Sigma\) -集的族，且对每个 \(i\in I\) 让 \(\varphi_i\) 是 \(\alpha\) 从E到 \(\mathbf{F}_i\) 。则映射 \((\varphi_i)_{i\in I}\) 将E映入 \(\prod_{i\in I}F_i\) （赋予乘积结构）是一个 \(\alpha\) -映射。

一个 \(\Sigma\) -集合 F 的子集 G 被称为 \(\Sigma\) -可容许的，如果 F 上的结构在 G 上诱导出一个物种结构 \(\Sigma\) （ \(\S\) 2，第 4 期）。

（CU \(_{III}\) 存在一个基数 \(\mathfrak{a}\) 具有以下性质：对于每一个 \(\Sigma\) -集合 \(\mathbf{F}\) 以及每一个 \(\alpha\) -映射 \(\varphi\) 的 \(\mathbf{E}\) 进入。 \(\mathbf{F}\) 存在一个 \(\Sigma\) -容许子集 \(\mathbf{G}\) 的 \(\mathbf{F}\) 它包含 \(\varphi(\mathbf{E})\) ，具有基数 \(\leqslant \mathfrak{a}\) ，使得 \(\mathbf{E}\) 进入。 \(\mathbf{G}\) 的映射具有与 \(\varphi\) 是 \(\alpha\) -映射，并且使得任意两个 \(\mathbf{G}\) 到 \(\Sigma\) -集合的态射，若它们在 \(\varphi(\mathbf{E})\) 是相等的。

CST22。如果条件 \((\mathrm{CU}_{\mathrm{I}})\) 至 \((\mathrm{CU}_{\mathrm{III}})\) 均成立，那么 E 的泛映射问题有解。

我们首先证明，如果存在一对 \((\mathbf{F}_{\mathbf{E}}, \varphi_{\mathbf{E}})\) 满足 \((\mathbf{AU}_{\mathbf{I}}^{\prime})\) ，那么对于 E 的泛映射问题也存在一个解。因为通过 \((\mathrm{CU}_{\mathrm{III}})\) 存在一个 \(\Sigma\) -容许子集 \(F_{E}^{\prime}\) 的 \(F_{E}\) 包含 \(\varphi_{E}(E)\) ，使得映射 \(\varphi_{E}^{\prime}\) 将E映入 \(F_{E}^{\prime}\) 其图与 \(\varphi_{E}\) 是 \(\alpha\) -映射，并且使得任意两个 \(F_{E}^{\prime}\) 到 \(\Sigma\) 在……上一致的集合 \(\varphi_{E}(E)\) 相等。设 j 为……的典范嵌入 \(F_{E}^{\prime}\) 进入。 \(F_{E}\) ，因此 \(\varphi_{E} = j \circ \varphi_{E}^{\prime}\) 。对于每一个态射 f，从 \(F_{E}\) 到 \(\Sigma\) -集合 F， \(f \circ j\) 是……的一个态射 \(F_{E}^{\prime}\) 到 F 中，并且我们有 \(f \circ \varphi_{E} = (f \circ j) \circ \varphi_{E}^{\prime}\) 。因此显然 \((\mathbf{F}_{\mathbf{E}}^{\prime}, \varphi_{\mathbf{E}}^{\prime})\) 满足 \((\mathbf{AU}_{\mathbf{I}}^{\prime})\) 并且 \((\mathbf{AU}_{\mathbf{II}}^{\prime})\) .

我们只需证明这样一个对子的存在性。 \((\mathrm{F}_{\mathbf{E}}, \varphi_{\mathbf{E}})\) 满足 \((\mathrm{AU}_{\mathbf{I}}^{\prime})\) . 设 \(s \in \mathrm{S}(x)\) 成为结构种类的典型刻画 \(\Sigma\) ，并考虑子集 L，它由 \(\mathfrak{P}(\mathfrak{a}) \times \mathrm{S}(\mathfrak{a}) \times \mathfrak{P}(\mathbf{E} \times \mathfrak{a})\) 中的所有三元组构成 \(\lambda = (\mathbf{X}, \mathbf{V}, \mathbf{P})\) 具有以下性质：“V是物种 \(\Sigma\) 在 \(X \subset a\) 上的一个结构，且 P 是某个 \(\alpha\) -映射（将E映射到X，关于结构V）的图''

（注意我们有 \(S(X) \subset S(a)\) ，这可以通过对阶梯构造方案的步骤长度逐步论证而容易看出 \(S\) ）。对于每个 \(\lambda = (X, V, P) \in L\) ，我们用 \(X_{\lambda}\) 集合 \(X\) 表示赋予结构 \(V\) 的集合，并用 \(\varphi_{\lambda}\) 表示 \(E\) 进入。 \(X_{\lambda}\) 其图形为 \(P\) .

让 \(F_{E}\) 是 \(\Sigma\) -集的映射，该映射是 \(X_{\lambda}\) 的乘积（它由 \((CU_{I})\) ），并令 \(\varphi_{E}\) 为映射 \(x \to (\varphi_{\lambda}(x))\) 将E映入 \(F_{E}\) ，这是一个 \(\alpha\) 根据 \((CU_{II})\) 的映射。让我们证明这一对 \((F_{E}, \varphi_{E})\) 满足 \((AU_{I}')\) 。给定一个从 E 到某个 \(\alpha\) -映射 \(\varphi\) -集合 F 的 \(\Sigma\) ，设 G 是 F 的一个子集，它满足 \((CU_{III})\) 中所述的条件。令 j 为 G 到 F 的典范嵌入，并令 \(\psi\) 为从 E 到 G 的映射，其图像与 \(\varphi\) ，因此 \(\varphi = j \circ \psi\) 相同。由此可知， \((CU_{III})\) 蕴含 \(\psi\) 是一个 \(\alpha\) -映射，从 E 到 G。由于 Card \((G) \leqslant a\) ，存在一个与 G 等势的子集 \(G'\) ，它包含于某个与 G 等势的集合中。设 g 为从 G 到 \(G'\) 如果我们通过 g 将物种结构 \(\Sigma\) 传输到 G 上，那么根据定义存在一个 \(\lambda\) L 中的元素，使得 \(G'\) （带有传输结构）等于 \(X_{\lambda}\) 并且使得 \(g \circ \psi = \varphi_{\lambda}\) 。然后 \(f = j \circ \overline{g}^1 \circ \mathrm{pr}_{\lambda}\) 是……的一个态射 \(\mathbf{F}_{\mathbf{E}}\) 进入。 \(\mathbf{F}\) 使得 \(\varphi = f \circ \varphi_{\mathbf{E}}\) ，证明完成。

CST23. 设 \((\mathbf{F}_{\mathbf{E}}, \varphi_{\mathbf{E}})\) 为 E 的泛映射问题的一个解。那么 \(\varphi_{E}\) 是 E 到 \(F_{E}\) 的单射，当且仅当对于 E 中每一对不同的元素 x, y，存在一个 \(\alpha\) -映射 \(\varphi\) -集合 F 的 \(\Sigma\) -集 F 使得 \(\varphi(x) \neq \varphi(y)\) .

由于 \(\varphi_{E}\) 是 \(\alpha\) -映射，该判据是定义的直接推论。

在这种情况下， \(\alpha\) -映射被称为分离E的元素，并且我们在术语上通常不区分E的元素与其在 \(\varphi_{E}\) 下的像。按照这一约定，如果 \((F_{E}, \varphi_{E})\) 是某个泛映射问题的解，并且如果条件 \((CU_{III})\) 被满足，那么每一个从E到某个 \(\alpha\) -集合F的 \(\Sigma\) -映射都唯一地延拓为从 \(F_{E}\) 到F的一个态射。

